% §8 结论
\label{sec:conclusion}

我们提出了因果相干性度量（$\branchCoherence$）——一个模型无关的、根植于操作系统原理的度量，用于量化从溯源图中提取的事件链的因果结构完整性——以及一个无监督校准协议，仅从良性数据中发现操作系统的自然因果视野。该度量的三分法逐对相干性函数区分身份保持桥接（进程）和有状态桥接（文件和套接字），根植于进程是独占能力而文件是共享资源这一OS原理。校准协议使用良性链上的肘部检测来设置所有参数，无需攻击标签或人工调优，并自动适配不同的OS配置。

在$\branchCoherence$指导下提取的因果链，由一个仅在良性事件序列上训练的自监督自回归Transformer进行验证，在DARPA E3上达到[X]\%的链级F1——超过基于窗口、基于随机游走、基于子图和基于后处理链重建的方法超过10个百分点的绝对F1——同时相比最先进的窗口级检测器减少超过50倍的调查成本。校准协议在E3和E5上发现不同的因果视野，展示了跨平台自动适配能力。

通过将APT检测从“异常在哪里/何时”重新定义为“攻击是哪个因果序列”，我们的方法为安全分析员提供了直接可操作的攻击叙事，而非可疑实体或时间窗口的列表。
